Pretende-se, futuramente, implementar \textit{scaling} vertical na infraestrutura do ESPEON, de modo a suportar turmas maiores e permitir a obtenção de amostras mais representativas. Há também interesse em publicar a extensão na \textit{Chrome Web Store}, possibilitando que docentes de todo o país utilizem a ferramenta de forma gratuita e descomplicada.

Além disso, considera-se a adoção de tecnologias para aprimoramento da interface dos \textit{pop-ups}, especialmente devido ao modelo atual de gerenciamento de estado das telas, realizado por referência no cliente. Por fim, vislumbra-se a criação de \textit{Bases do Conhecimento} para ingestão em \textit{LLMs}, permitindo fornecer aos docentes \textit{insights} sobre suas aulas por meio de diagnósticos em linguagem natural, contribuindo para processos de reflexão e aprimoramento contínuo no aprendizado.